\uuid{yqxX}
\exo7id{3585}
\titre{exo7 3585}
\auteur{quercia}
\organisation{exo7}
\datecreate{2010-03-10}
\isIndication{false}
\isCorrection{true}
\chapitre{Réduction d'endomorphisme, polynôme annulateur}
\sousChapitre{Applications}
\cpge{mpsi-mp2i.calcul-matriciel-systemes.e, pcsi.determinants.c, ptsi.determinants.a, mpsi-mp2i.applications-lineaires.b, pcsi.applications-lineaires.b, ptsi.applications-lineaires.b, bcpst1.matrices, mp-mpi.reduction-endomorphismes.b, pc.reduction.a, psi.reduction-endomorphismes.a, pt.reduction-endomorphismes.a, mp-mpi.reduction-endomorphismes.c, pc.reduction.b, psi.reduction-endomorphismes.b, pt.reduction-endomorphismes.b, mp-mpi.reduction-endomorphismes.d, pc.reduction.c, psi.reduction-endomorphismes.c, pt.reduction-endomorphismes.c, bcpst2.valeurs-propres-vecteurs-propres.a, bcpst2.valeurs-propres-vecteurs-propres.b}
\module{Algèbre}
\niveau{L2}
\difficulte{}

\contenu{
\texte{
Soit $A=\begin{pmatrix}-1&2&1\cr2&-1&-1\cr-4&4&3\cr\end{pmatrix}$.
}
\begin{enumerate}
    \item \question{Calculer $A^n$.}
    \item \question{Soit $U_0=\begin{pmatrix}-2\cr4\cr1\end{pmatrix}$ et $(U_n)$ défini par la relation~:
    $U_{n+1} = AU_n$. Calculer $U_n$ en fonction de~$n$.}
    \item \question{Soit $X(t) = \begin{pmatrix}x(t)\cr y(t)\cr z(t)\cr\end{pmatrix}$.
    Résoudre $\frac{d X}{d t} = AX$.}
\reponse{
$A^{2k}=I$, $A^{2k+1}=A$.
}
\end{enumerate}
}
